\begin{abstract}
A picture enters a vision-language model as thousands of tokens and leaves it as a few sentences of text.
The pages of a comic multiply that input while the answer stays a single paragraph.
Speculative decoding lets a small draft model propose several tokens that the target model verifies in one pass, and recent drafters for vision-language models compress the image or hide it from the draft model entirely.
We evaluate the drafters \vispec and \hivis against autoregressive decoding on single-image question answering and captioning and on story summarization over windows of manga pages from \rvp.
We also prune the draft model's visual input and evict image entries from the target's cache.
An exactness protocol traces every disagreement with autoregressive decoding to a near-tie in the target's own logits, and full precision reproduces the autoregressive output on every sample it replays.
Speculative decoding roughly halves the time of a long answer.
Yet a draft model that never sees the image is as fast as one that does, and evicting most image entries leaves the cost of a decoding step almost unchanged.
The cost of the image falls before the first token, where no drafter acts.
Its cache entries are a few percent of the bytes each decoding step reads.
On manga pages prompt processing takes close to half of a speculative decoder's time and holds its end-to-end gain below two times.
\end{abstract}
